% Options for packages loaded elsewhere
\PassOptionsToPackage{unicode}{hyperref}
\PassOptionsToPackage{hyphens}{url}
\documentclass[
]{article}
\usepackage{xcolor}
\usepackage{amsmath,amssymb}
\usepackage{amsthm}
\usepackage{booktabs}
\setcounter{secnumdepth}{3}
\usepackage{iftex}
\ifPDFTeX
  \usepackage[T1]{fontenc}
  \usepackage[utf8]{inputenc}
  \usepackage{textcomp}
\else
  \usepackage{unicode-math}
  \defaultfontfeatures{Scale=MatchLowercase}
  \defaultfontfeatures[\rmfamily]{Ligatures=TeX,Scale=1}
\fi
\usepackage{lmodern}
\IfFileExists{upquote.sty}{\usepackage{upquote}}{}
\IfFileExists{microtype.sty}{%
  \usepackage[]{microtype}
  \UseMicrotypeSet[protrusion]{basicmath}
}{}
\makeatletter
\@ifundefined{KOMAClassName}{%
  \IfFileExists{parskip.sty}{%
    \usepackage{parskip}
  }{% else
    \setlength{\parindent}{0pt}
    \setlength{\parskip}{6pt plus 2pt minus 1pt}}
}{%
  \KOMAoptions{parskip=half}}
\makeatother
\setlength{\emergencystretch}{3em}
\providecommand{\tightlist}{%
  \setlength{\itemsep}{0pt}\setlength{\parskip}{0pt}}
\usepackage{bookmark}
\IfFileExists{xurl.sty}{\usepackage{xurl}}{}
\urlstyle{same}
\hypersetup{
  pdftitle={The Trinity as a Consequence of a Monistic Modal Ontology under Superposition Axioms – A Metalogical Investigation},
  hidelinks,
  pdfcreator={LaTeX via pandoc}
}

\newtheorem{theorem}{Theorem}
\newtheorem{lemma}{Lemma}
\newtheorem{definition}{Definition}
\newtheorem{corollary}{Corollary}

\title{The Trinity as a Consequence of a Monistic Modal Ontology under Superposition Axioms \\ \large A Metalogical Investigation}
\author{Paul Koop}
\date{}

\begin{document}
\maketitle

\begin{center}
\emph{Past and future are horizons of knowledge. \\ The present is the locus of reality. \\ Nothing that was, is lost.}
\end{center}

\newpage
{
\setcounter{tocdepth}{3}
\tableofcontents
}
\newpage

\section{Preface: What This Work Is – And What It Is Not}

This work explicitly does not aim to prove the truth of a metaphysical or theological claim.

Rather, it investigates the \textbf{logical structure of an axiomatic system}. The question is not whether the Trinitarian structure exists, but \textbf{which additional assumptions are required for it to follow from a monistic modal ontology}.

The investigation is therefore \textbf{metalogical}: it determines a sufficient metaphysical supplementary structure – not physical mechanisms or empirical facts. The question of the minimality of this structure remains open and is a subject for future research.

\subsection{Four Levels of Investigation}

The work distinguishes four levels:

\begin{table}[h]
\centering
\begin{tabular}{lll}
\toprule
\textbf{Level} & \textbf{Question} & \textbf{Method} \\
\midrule
\textbf{Logic} & Which structures are consistent? & Formal proofs, tableaux, model theory. \\
\textbf{Metaphysics} & Which structural conditions are sufficient? & Axiomatic analysis. \\
\textbf{Physics} & Through which processes are these conditions realized? & Empirical science, technology. \\
\textbf{Theology/Hermeneutics} & How are these structures interpreted? & Interpretation, linguistic usage, tradition. \\
\bottomrule
\end{tabular}
\end{table}

The formal core investigation operates \textbf{exclusively} on the first two levels. The later parts (III and the hermeneutical outlook) show how these formal structures can connect to physical and hermeneutical questions without answering them themselves.

\subsection{On the Use of the Term "Superposition"}

The term "superposition" is used in this work \textbf{axiomatically}. It is not claimed that physical quantum mechanics already proves this ontological role. The superposition axioms ASP1–ASP5 are \textbf{metaphysical assumptions}, not physical theorems. They formalize the idea of a primordial state in which all possibilities are undividedly contained – independent of whether and how this state is physically realized.

The central claim of this work is:

\begin{quote}
\textbf{The investigation determines a sufficient metaphysical supplementary structure from which the Trinitarian structure follows. The question of its minimality remains open.}
\end{quote}

\section{Criteria for the Choice of Axioms}

The axioms used in this work are not arbitrary. They formalize certain philosophical intuitions rooted in the tradition of monistic metaphysics:

\begin{itemize}
\item \textbf{A1 (Monism):} The unity of reality – a fundamental intuition of Parmenides, Neoplatonism, and the idealist tradition.
\item \textbf{A11 (Transcendental Bridge):} The impossibility of an unknowable world – a Kantian insight stating that the conditions of experience are also the conditions of the objects of experience.
\item \textbf{A12 (Experienceability ↔ Realization):} The identity of experienceability and realization – a pragmatist or transcendental assumption.
\item \textbf{ASP1–ASP5 (Superposition):} The idea of a primordial state in which all possibilities are undividedly contained – an intuition found in speculative metaphysics (e.g., in Schelling or Neoplatonism), but also serving as an analogy in quantum physics.
\end{itemize}

The choice of axioms is guided by the goal of finding a \textbf{sufficient} structure from which the Trinitarian structure follows from a monistic modal ontology. Other axiomatizations are possible and could lead to different consequences. The axioms are therefore not to be understood as ultimate truths, but as \textbf{hypothetical assumptions} whose consequences are investigated. This corresponds to the methodological self-understanding of the work as a metalogical investigation.

\section{Introduction: Methodological Self-Restriction}

Classical metaphysics has repeatedly attempted to derive the fundamental structure of reality from a few basic principles. A particular challenge arises from three seemingly distinct aspects:

\begin{enumerate}
\item \textbf{Why is there something rather than nothing?} – origin (U).
\item \textbf{How does the totality of all possibilities relate to actuality?} – totality (T).
\item \textbf{How can a state arise within reality that recognizes reality itself?} – self-knowledge (S).
\end{enumerate}

The ontology examined here proposes to answer these three questions not through three separate metaphysical substances, but through three necessary perspectives of the same reality.

The present work, however, is \textbf{not a proof of the existence} of this structure. It is a \textbf{metalogical investigation} into the axioms needed for this structure to follow from a monistic modal ontology.

\section{Formal Language and Logical Framework}

\subsection{Modal Logic S5}

We work in first-order modal logic with identity in system S5:

\begin{itemize}
\item \textbf{Necessity:} \(\Box p\)
\item \textbf{Possibility:} \(\Diamond p := \neg\Box\neg p\)
\end{itemize}

\textbf{Axioms of S5:}
\begin{itemize}
\item (K) \(\Box(p \rightarrow q) \rightarrow (\Box p \rightarrow \Box q)\)
\item (T) \(\Box p \rightarrow p\)
\item (4) \(\Box p \rightarrow \Box\Box p\)
\item (5) \(\Diamond p \rightarrow \Box\Diamond p\)
\end{itemize}

\textbf{Inference Rules:}
\begin{itemize}
\item \textbf{Modus Ponens:} From \(p\) and \(p \rightarrow q\), infer \(q\).
\item \textbf{Necessitation:} From \(p\) (derivable), infer \(\Box p\).
\end{itemize}

\textbf{Tableau Rules for S5:}
\begin{align*}
& (\neg\Box) & \frac{\neg\Box A}{\Diamond\neg A} \\
& (\Diamond) & \frac{\Diamond A}{A @ w_{\text{new}}} \quad \text{(new world)} \\
& (\Box) & \frac{\Box A @ w}{A @ v} \quad \text{for every already existing world } v \\
& (\neg\forall) & \frac{\neg\forall x A}{\exists x \neg A} \\
& (\exists) & \frac{\exists x A}{A[a/x]} \quad \text{(a new)} \\
& (\forall) & \frac{\forall x A}{A[a/x]} \quad \text{(a arbitrary)}
\end{align*}

\subsection{Predicate Logic}

We use classical first-order predicate logic with identity (\(=\)) and the usual quantifiers (\(\forall, \exists\)).

\section{The Axioms (Basic)}

\subsection{A1 – Monism}

There are no fundamentally separated domains of reality.

\[
\boxed{A1 := \Box\neg\exists x\exists y\, FundamentalSeparated(x,y)}
\]

\subsection{A4 – Existence of a Possible Reality}

There is at least one possible world.

\[
\boxed{A4 := \Diamond\exists w\, World(w)}
\]

\subsection{A11 – Transcendental Bridge}

A world is separated from consciousness exactly when it contains no consciousness.

\[
\boxed{A11 := \forall w \left( \text{WorldSeparatedFromConsciousness}(w) \leftrightarrow \neg \exists c (Consciousness(c) \land c(w)) \right)}
\]

\subsection{A12 – Experienceability and Realization}

\[
\boxed{A12 := \forall p. \text{Experienceable}(p) \leftrightarrow \exists w. Realized(w, p)}
\]

\subsection{A13, A14, A15 – The Three Implications of the Cases}

\[
\boxed{A13 := \Box\forall x(T(x) \rightarrow U(x))}
\]
\[
\boxed{A14 := \Box\forall x(U(x) \rightarrow S(x))}
\]
\[
\boxed{A15 := \Box\forall x(S(x) \rightarrow T(x))}
\]

\section{Definition of the Limit Structure}

\subsection{Totality T as an Open, Well-Founded Interval}

\[
\boxed{T := \{ x \mid U < x < S \}}
\]

\subsection{Origin U as Lower Limit}

\[
\boxed{U := \lim_{x \to \inf} T}
\]

\subsection{Self-Knowledge S as Upper Limit}

\[
\boxed{S := \lim_{x \to \sup} T}
\]

\subsection{The Trinity}

\[
\boxed{Tr := U \land T \land S}
\]

\section{PART I: THE METALOGICAL CORE – NON-DERIVABILITY IN PURE S5}

\subsection{Goal}

Examine whether:

\[
\boxed{\text{S5} \;\vdash\; \Box\forall x\,Tr(x)}
\]

follows from the axioms \(A1, A4, A11, A12, A13, A14, A15\).

\subsection{Method}

Construct an \textbf{open S5 tableau} for the negation of the target formula.  
According to the \textbf{soundness and completeness theorem} of the S5 tableau calculus:

\begin{quote}
A set of formulas is \textbf{satisfiable} in an S5 model exactly when the tableau has \textbf{an open branch}.
\end{quote}

\subsection{Result}

The tableau has an open branch. Therefore:

\begin{theorem}[Non-derivability in pure S5]
\[
\boxed{
\text{S5} \;\not\vdash\; \Box\forall x\,Tr(x)
}
\]
\end{theorem}

\textbf{Interpretation:} The axioms \(A1, A4, A11, A12, A13, A14, A15\) are \textbf{not sufficient} to force the Trinitarian structure.

This is an \textbf{independence result} in the sense of the metalogical tradition: It shows that the Trinitarian structure does not follow from the given axioms – regardless of whether it is true. The actual result is: \textbf{Without additional structure, it does not work.}

\section{PART II: A SUFFICIENT EXTENSION – CONSTRUCTION AND PROOF}

The superposition axioms are not formulated independently of the target structure, but constitute a \textbf{constructive extension} of the original ontology. They are chosen so as to produce the desired consequence. Their constructive nature means: They are a \textbf{means} to produce the desired consequence – not a discovery of independent facts.

\subsection{The Superposition Extension (ASP)}

\[
\begin{aligned}
ASP1 &:= C @ w_{\text{super}} \\
ASP2 &:= \forall v (w_{\text{super}} R v \rightarrow (C @ v \leftrightarrow v = w_{\text{super}})) \\
ASP3 &:= \forall x\,Tr(x) @ w_{\text{super}} \\
ASP4 &:= \forall v (w_{\text{super}} R v \rightarrow \forall x\,Tr(x) @ v) \\
ASP5 &:= \forall v (v \neq w_{\text{super}} \rightarrow w_{\text{super}} R v)
\end{aligned}
\]

\subsection{Model-Theoretic Formulation}

The class of all S5 models \textbf{whose interpretation additionally satisfies the axioms ASP1–ASP5} also satisfies the Trinitarian formula \(\Box\forall x\,Tr(x)\).

\subsection{Proof of Derivability in S5+SP}

\begin{theorem}[Derivability in S5+SP]
\[
\boxed{
\text{S5+SP} \;\vdash\; \Box\forall x\,Tr(x)
}
\]
\end{theorem}

\textbf{Interpretation:} The superposition axioms ASP1–ASP5 are \textbf{sufficient} to derive the Trinity. Whether they are necessary – whether each individual axiom is indispensable – is a subject for future minimality investigations.

\section{PART III: ESCHATOLOGICAL CONSEQUENCES – A FURTHER SUFFICIENT EXTENSION}

The statements about resurrection, unitarity, and integration are considerably more dependent on the introduced axioms than the Trinitarian structure itself. Their persuasiveness therefore depends less on the formal derivation than on the \textbf{plausibility of the axioms}. The work offers no independent justification here.

\subsection{The Extension of S}

\[
\boxed{A16 := \forall c(\text{Consciousness}(c) \rightarrow \text{integrated}(c, S))}
\]

\[
\boxed{A17 := \forall c(\text{Knowledge}(c, S) \leftrightarrow \text{integrated}(c, S))}
\]

\subsection{Unitarity}

\[
\boxed{A18 := \Box\forall I(\text{Existence}(I) \rightarrow \text{Conservation}(I))}
\]

\subsection{Death as Transition}

\[
\boxed{A19 := \forall c(\text{Death}(c) \rightarrow \text{Integration}(c, S))}
\]

\[
\boxed{A20 := \forall c(\text{Knowledge}(c, S) \leftrightarrow \text{Existence}(c))}
\]

\subsection{Theorem: Resurrection}

\[
\boxed{\text{S5+SP+U} \;\vdash\; \Box\forall c(\text{Consciousness}(c) \rightarrow \text{Resurrection}(c))}
\]

\textbf{Interpretation:} The axioms A16–A20 are \textbf{sufficient} to derive resurrection as a theorem. Whether they are minimal or whether individual axioms are dispensable remains open.

\section{Hermeneutical Outlook: Three Movements of Thought as a Resonance Space}

The following considerations are \textbf{not part of the formal proof}, but a hermeneutical outlook showing the scope of the formal structure in other contexts. The reader here consciously leaves the level of formal logic and enters the space of interpretation.

The formal structure developed in this work finds its existential and poetic correspondence in three independent movements of thought that can serve as a resonance space for the reflections presented here.

\subsection{The Unitarity of Pain}

The first movement of thought shows that suffering is not meaningless, but the price of differentiation necessary for knowledge. The unitarity of the universe guarantees that nothing is lost – not even pain.

\begin{quote}
\textbf{The black hole of death is transformed through resurrection into the Hawking radiation of meaning.}
\end{quote}

\subsection{Inverse Christology according to Evolution}

The second movement of thought shows that Christ is not the metaphysical starting point, but the historical culmination point of evolution. The incarnation is the condensation of the cosmic process.

\begin{quote}
\textbf{The incarnation is the local manifestation of superposition. The resurrection is its global integration into S.}
\end{quote}

\subsection{The Omega Poems}

The third movement of thought is a lyrical approach to the limit. It shows that the formal structure of the Trinity is not merely abstract, but comes alive in the language of poetry.

\begin{quote}
\textbf{The poetry of the Omega Point is the language of omniscience, uniting all consciousnesses within itself.}
\end{quote}

\section{Conclusion: What Has Been Shown?}

The present investigation has shown:

\begin{enumerate}
\item The Trinitarian structure is \textbf{not derivable in pure S5} (Part I). This is the actual metalogical finding – an independence result.
\item The superposition axioms ASP1–ASP5 are \textbf{sufficient} to derive the Trinity (Part II). The question of their minimality remains open.
\item The axioms A16–A20 are \textbf{sufficient} to derive unitarity and resurrection as theorems (Part III). Here too, the question of minimality remains open.
\end{enumerate}

\begin{quote}
\textbf{The central claim of this work is not: "The Trinity exists." Rather: "The investigation determines a sufficient metaphysical supplementary structure from which the Trinitarian structure follows. The question of its minimality remains open."}
\end{quote}

The work is thus not an ontological claim, but a \textbf{metalogical investigation into the dependence of a structure on an axiomatic system}. It shows a sufficient extension of the axiomatic system under which the Trinitarian structure follows. Whether this extension is minimal or unique remains open.

The work shows that the Trinitarian structure is not an arbitrary metaphysical claim, but a \textbf{consequence} of certain axiomatic decisions. It thus makes the structure \textbf{verifiable} and \textbf{debatable} – independent of theological presuppositions. The question of whether the axioms proposed here are minimal or whether they can be reduced remains open for future investigations. Equally open remains the question of how these structures could be physically realized or theologically interpreted – this is a matter for physics and hermeneutics, not logic.

\begin{center}
\emph{This treatise was written in the spirit of rigorous modal logic, yet in the language of philosophy – for truth requires both: the precision of the formula and the breadth of the concept.}
\end{center}

\end{document}